\documentclass[11pt,a4paper]{article}
\usepackage[margin=1.05in]{geometry}
\usepackage{amsmath,amssymb,amsthm,mathtools}
\usepackage{lmodern}
\usepackage[T1]{fontenc}
\usepackage{microtype}
\usepackage{booktabs}
\usepackage{longtable}
\usepackage{xurl}
% Lean 4 listings: Unicode source glyphs are mapped to LaTeX math, so no
% Unicode font is required (Unicode-safe under pdfLaTeX, XeTeX and tectonic).
\usepackage{listings}
\usepackage{xcolor}
\usepackage{newunicodechar}
\definecolor{leankw}{RGB}{0,64,128}
\definecolor{leancomment}{RGB}{96,96,96}
\lstdefinelanguage{lean4}{
  morekeywords={theorem,lemma,def,structure,fun,by,have,show,from,intro,exact,
    refine,simp,simpa,apply,namespace,end,open,variable,noncomputable,Prop,Type,rfl},
  sensitive=true, morecomment=[l]{--}, morecomment=[s]{/-}{-/},
  morestring=[b]"}
\lstset{language=lean4, basicstyle=\ttfamily\small, keywordstyle=\color{leankw}\bfseries,
  commentstyle=\color{leancomment}\itshape, columns=fullflexible, keepspaces=true,
  breaklines=true, breakatwhitespace=true, frame=tb, framerule=0.4pt,
  aboveskip=0.8\baselineskip, belowskip=0.8\baselineskip, xleftmargin=0.5em,
  mathescape=true}

% Unicode Lean glyphs rendered as LaTeX math (works inside and outside listings)
\newunicodechar{ℝ}{\ensuremath{\mathbb{R}}}
\newunicodechar{ℕ}{\ensuremath{\mathbb{N}}}
\newunicodechar{ℤ}{\ensuremath{\mathbb{Z}}}
\newunicodechar{ℂ}{\ensuremath{\mathbb{C}}}
\newunicodechar{→}{\ensuremath{\to}}
\newunicodechar{←}{\ensuremath{\leftarrow}}
\newunicodechar{↔}{\ensuremath{\leftrightarrow}}
\newunicodechar{∀}{\ensuremath{\forall}}
\newunicodechar{∃}{\ensuremath{\exists}}
\newunicodechar{≤}{\ensuremath{\le}}
\newunicodechar{≥}{\ensuremath{\ge}}
\newunicodechar{≠}{\ensuremath{\ne}}
\newunicodechar{¬}{\ensuremath{\neg}}
\newunicodechar{∧}{\ensuremath{\wedge}}
\newunicodechar{∨}{\ensuremath{\vee}}
\newunicodechar{∈}{\ensuremath{\in}}
\newunicodechar{∫}{\ensuremath{\int}}
\newunicodechar{⁻}{\ensuremath{^{-}}}
\newunicodechar{¹}{\ensuremath{^{1}}}
\newunicodechar{²}{\ensuremath{^{2}}}
\newunicodechar{⟨}{\ensuremath{\langle}}
\newunicodechar{⟩}{\ensuremath{\rangle}}
\newunicodechar{‖}{\ensuremath{\|}}
\newunicodechar{∘}{\ensuremath{\circ}}
\newunicodechar{•}{\ensuremath{\bullet}}
\newunicodechar{·}{\ensuremath{\cdot}}
\newunicodechar{⬝}{\ensuremath{\cdot}}
\newunicodechar{∑}{\ensuremath{\sum}}
\newunicodechar{λ}{\ensuremath{\lambda}}
\newunicodechar{κ}{\ensuremath{\kappa}}
\newunicodechar{ω}{\ensuremath{\omega}}
\newunicodechar{ν}{\ensuremath{\nu}}
\newunicodechar{π}{\ensuremath{\pi}}
\newunicodechar{α}{\ensuremath{\alpha}}
\newunicodechar{β}{\ensuremath{\beta}}
\newunicodechar{ε}{\ensuremath{\varepsilon}}
\newunicodechar{δ}{\ensuremath{\delta}}
\newunicodechar{∇}{\ensuremath{\nabla}}
\newunicodechar{∂}{\ensuremath{\partial}}
\newunicodechar{⊥}{\ensuremath{\perp}}
\newunicodechar{⊤}{\ensuremath{\top}}
\newunicodechar{₀}{\ensuremath{_0}}
\newunicodechar{₁}{\ensuremath{_1}}
\newunicodechar{₂}{\ensuremath{_2}}
\newunicodechar{ᵀ}{\ensuremath{^T}}
\newunicodechar{≫}{\ensuremath{\gg}}
\newunicodechar{⦃}{\ensuremath{\{\!|}}
\newunicodechar{⦄}{\ensuremath{|\!\}}}
\newunicodechar{∞}{\ensuremath{\infty}}
\newunicodechar{×}{\ensuremath{\times}}
\newunicodechar{⊆}{\ensuremath{\subseteq}}
\newunicodechar{≈}{\ensuremath{\approx}}
\newunicodechar{𝔗}{\ensuremath{\mathfrak{T}}}
\newunicodechar{⇑}{\ensuremath{\Uparrow}}
\newunicodechar{ᶜ}{\ensuremath{^{c}}}

% Additional glyphs used by the spatial-bridge paper listings.
\newunicodechar{₊}{\ensuremath{_{+}}}
\newunicodechar{ˢ}{\ensuremath{^{s}}}
\newunicodechar{ξ}{\ensuremath{\xi}}
\newunicodechar{γ}{\ensuremath{\gamma}}
\newunicodechar{Ω}{\ensuremath{\Omega}}

\usepackage[colorlinks,linkcolor=leankw,citecolor=leankw,urlcolor=leankw]{hyperref}

\theoremstyle{plain}
\newtheorem{theorem}{Theorem}[section]
\newtheorem{lemma}[theorem]{Lemma}
\newtheorem{proposition}[theorem]{Proposition}
\newtheorem{corollary}[theorem]{Corollary}
\theoremstyle{definition}
\newtheorem{definition}[theorem]{Definition}
\newtheorem{remark}[theorem]{Remark}

\newcommand{\lean}[1]{{\upshape\path{#1}}}
\newcommand{\leanu}[1]{\textnormal{\texttt{#1}}}
% \declref{declaration}{repository file}{line}: machine-checked by check_refs.py
\newcommand{\declref}[3]{\textnormal{\lean{#1} (\lean{#2}:#3)}}
\newenvironment{leanblock}{\par\smallskip\noindent\rule{\linewidth}{0.4pt}\par\noindent\begin{minipage}{\linewidth}\raggedright\ttfamily\small}{\end{minipage}\par\noindent\rule{\linewidth}{0.4pt}\par\smallskip}
\newcommand{\kw}[1]{\textcolor{leankw}{\textbf{#1}}}
\Urlmuskip=0mu\relax
\newcommand{\R}{\mathbb{R}}
\newcommand{\T}{\mathbb{T}}
\newcommand{\Ennr}{[0,\infty]}
\DeclareMathOperator{\curl}{curl}
\renewcommand{\div}{\operatorname{div}}
\emergencystretch=3em
\sloppy

\newcommand{\mainsha}{43653aab3b51269a71a713ccbc4d892869bbe497}
\newcommand{\proofsha}{60e7ed0}

\title{\textbf{A Kernel-Checked Spatial Bridge for Vortex-Line Comparison:\\
Lean~4 Formalization of the Deng--Hou--Yu Spatial Identity\\
for Smooth Periodic Navier--Stokes Fields}}
\author{Frohman, Benjamin Stanley\\
\small ORCID: \href{https://orcid.org/0009-0006-7068-3718}{0009-0006-7068-3718}\quad|\quad
\href{mailto:frohmanbenjamin@gmail.com}{\texttt{frohmanbenjamin@gmail.com}}\\
\small Repository: \href{https://github.com/BenFrohman/Clay-}{\lean{https://github.com/BenFrohman/Clay-}}}
\date{November 9, 2026 (America/Chicago, UTC$-6$). Preprint; not peer reviewed.}

\begin{document}
\maketitle

\begin{abstract}
We describe a Lean~4/Mathlib formalization, merged in the public repository
\lean{BenFrohman/Clay-} at commit \texttt{43653aa}, of the \emph{spatial} half
of Theorem~1 of Deng, Hou and Yu (DHY): along a nonvanishing vortex-line
segment from a point $x$ to a point $y$, the vorticity magnitudes satisfy
$|\omega(x)|\le e^{C}|\omega(y)|$ whenever
$\bigl|\int_0^{L}(\div\xi)(\gamma(s))\,ds\bigr|\le C$, where
$\xi=\omega/|\omega|$. The library proves the named target
\lean{SpatialBridgeObligation} for every packet of solution-linked
vortex-line data over every smooth periodic Navier--Stokes solution in the
repository's encoding, with no assumed comparison, regularity supplement or
integrability premise. The proof derives spatial $C^2$ regularity from the
solution hypothesis, mixed-partial symmetry, $\div\curl=0$, a product-rule
identity for $|\omega|$ and $\xi$, a logarithmic chain rule along the curve,
and the fundamental theorem of calculus. As a consequence the library proves
that finiteness of the time integral of the endpoint vorticity implies
finiteness of the Beale--Kato--Majda integral. All of the theorems are checked
by the Lean kernel, and the full library axiom audit (125 theorems, 236
declarations) reports only \texttt{propext}, \texttt{Classical.choice} and
\texttt{Quot.sound}, with no \texttt{sorry}. We write every step of the
argument out on paper. We also record three limitations, each of them
kernel-checked: the argument uses only spatial $C^2$ smoothness, so it is
kinematic and not specific to Navier--Stokes; the endpoint hypothesis is
\emph{equivalent} to Beale--Kato--Majda finiteness, so the criterion moves the
analytic difficulty rather than removing it; and at the pinned Mathlib, the
repository's solution predicate requires joint real-analyticity rather than
$C^\infty$ smoothness. An informal, unformalized shear-flow argument shows
that, granting the classical uniqueness of smooth solutions, the
repository's formal version of Statement~(B) is false as written. That formal
statement should be corrected; doing so does not affect the results here. \textbf{This is not a proof of the Clay Millennium
Navier--Stokes problem, and it does not prove global regularity for any class
of data.} Kernel verification certifies the Lean statements as written. It
does not certify their fidelity to the intended PDE notions.
\end{abstract}

\tableofcontents

%==========================================================================
\section{Scope and non-claims}\label{sec:scope}

\paragraph{What this paper is.} It is an exposition of a small, completely
formalized piece of classical vorticity kinematics. It gives the exact Lean
statements, file and line references at a pinned commit, a complete
human-readable proof that follows the Lean proof step by step, and
reproduction instructions. The mathematics is classical: it is the identity
(3.1)--(3.2) of DHY Lemma~1 and the ratio bound (2.2) of DHY
Theorem~1~\cite{DHY2005}. \emph{No novelty is claimed for the mathematics.}
The contribution is the machine-checked formalization, its tie to the
repository's definition of a Navier--Stokes solution, and a precise account of
what remains open.

\paragraph{What this paper is not.}
\begin{enumerate}
  \item It is \textbf{not} a proof or a disproof of Fefferman's Statement~(B)~\cite{Fefferman2006} or of any other Clay Millennium statement.
  \item It does \textbf{not} prove that any nonzero Navier--Stokes solution admits a vortex-line packet (Definition~\ref{def:packet}). It also does not prove that a packet with time-uniform constants and measurable dependence on time exists.
  \item It does \textbf{not} prove the endpoint integrability hypothesis (Definition~\ref{def:endpoint}) for any nonzero solution. Proposition~\ref{prop:equiv} shows that this hypothesis is exactly as strong as finiteness of the Beale--Kato--Majda (BKM) integral.
  \item It does \textbf{not} formalize the classical local existence theory or the BKM continuation theorem~\cite{BKM1984}. In the repository these appear only as the named, unproved proposition \lean{PeriodicMaximalSolutionWithBKM}.
  \item It does \textbf{not} formalize DHY Theorem~2, which is specific to the Euler equations, or the degenerate and general cases covered by DHY Theorem~1.
  \item A successful kernel check means that the Lean terms typecheck against the Lean definitions. It does \textbf{not} show that those definitions faithfully encode Fefferman's notions of a smooth periodic solution, or any physical statement. Section~\ref{sec:analytic} records a concrete fidelity issue of exactly this kind.
\end{enumerate}

%==========================================================================
\section{The formal setting}\label{sec:setting}

All declarations live in the repository~\cite{ClayRepo} at commit
\texttt{\mainsha} on \texttt{main}. That commit is the merge of pull
request~\#9, whose proof commit is \texttt{\proofsha}. File and line numbers
below refer to that commit. In references of the form
\emph{name (file:line)}, the file is relative to \lean{lean/Clay/}; the pull request carrying this paper does not
change any Lean library file. Lean is version~4.16.0 and Mathlib is tag
\texttt{v4.16.0} (revision
\texttt{a6276f4c6097675b1cf5ebd49b1146b735f38c02})~\cite{Lean4,Mathlib2020}.

\subsection{Space, derivatives and vorticity}
Let $E=\R^3$ with the Euclidean norm (\leanu{E3 := EuclideanSpace ℝ (Fin 3)},
\lean{lean/Clay/StatementB.lean}:35), and let $e_1,e_2,e_3$ be the standard
basis (\declref{e}{StatementB.lean}{38}). In Lean the indices are
$0,1,2$; we write $1,2,3$. For $f\colon E\to\R$ the partial derivative is
defined through the Fr\'echet derivative,
\[
  \partial_j f(x) := Df(x)[e_j]
\]
(\declref{pd}{StatementB.lean}{41}).
Mathlib's \lean{fderiv} returns $0$ at points where $f$ is not differentiable.
Every use of $\partial_j$ below is at a point where differentiability has been
proved, so this convention never enters. For $v\colon E\to E$,
\[
  \curl v=(\partial_2v_3-\partial_3v_2,\ \partial_3v_1-\partial_1v_3,\ \partial_1v_2-\partial_2v_1)
\]
(\declref{curl}{StatementB.lean}{53}),
and $\div w=\sum_i\partial_iw_i$; the predicate $\div w\equiv 0$ is
\declref{DivFree}{StatementB.lean}{50}.

\subsection{Solutions, lifespans and the BKM integral}
For $T\in\Ennr$ the lifespan is $[0,T)=\{t:0\le t,\ \mathrm{ofReal}(t)<T\}$
(\declref{lifespanSet}{StatementB.lean}{82}). The predicate
\declref{NSOn}{StatementB.lean}{65}, written
$\mathrm{NSOn}(\nu,S,u_0,u,p)$, asserts the following.
\begin{enumerate}
  \item[(N1)] $(t,x)\mapsto u(t,x)$ and $(t,x)\mapsto p(t,x)$ satisfy \leanu{ContDiffOn ℝ ⊤} on $S\times\R^3$.
  \item[(N2)] The momentum equation holds with the one-sided time derivative \lean{derivWithin} along $S$.
  \item[(N3)] $\div u(t)=0$ for $t\in S$.
  \item[(N4)] $u(0,\cdot)=u_0$.
  \item[(N5)] $u(t)$ and $p(t)$ are $1$-periodic in each coordinate for $t\in S$.
\end{enumerate}
The vorticity supremum is
$\Omega(t)=\sup_x|\curl u(t)(x)|\in\Ennr$
(\declref{vortSup}{StatementB.lean}{95}), and
\[
  \mathrm{BKMIntegralFinite}(u,T)\ :\iff\ \int^{-}_{[0,T)}\Omega(t)\,dt<\infty
\]
(\declref{BKMIntegralFinite}{StatementB.lean}{98}),
where $\int^{-}$ is the lower Lebesgue integral of a $\Ennr$-valued function.
The quantifier \declref{ForAllFiniteSolutions}{StatementB.lean}{102}
ranges over all $\nu>0$, all data $u_0$ satisfying \lean{SmoothPeriodicData},
all $(u,p)$, and all finite $T$ with $\mathrm{NSOn}(\nu,[0,T),u_0,u,p)$.

\subsection{A fidelity remark: \texorpdfstring{\leanu{⊤}}{top} means analytic at the pinned Mathlib}\label{sec:analytic}
At the pinned Mathlib the smoothness index of \lean{ContDiff} and
\lean{ContDiffOn} has type \leanu{WithTop ℕ∞}. Mathlib writes $\infty$ for
$((\top:\mathbb{N}_\infty):\leanu{WithTop ℕ∞})$, which means $C^\infty$, and
$\omega$ for $\top:\leanu{WithTop ℕ∞}$, which means real-analytic (Mathlib
file \lean{Mathlib/Analysis/Calculus/ContDiff/Defs.lean}, ``Notations'').
\lean{NSOn} and \lean{SmoothPeriodicData} both write a bare \leanu{⊤}, so the
index is $\omega$. The paper-check file
\lean{papers/spatial-bridge/PaperChecks.lean} proves this explicitly.
\begin{itemize}
  \item \lean{NSOn_velocity_analyticOn}: $\mathrm{NSOn}(\nu,S,u_0,u,p)$ implies that $(t,x)\mapsto u(t,x)$ is real-analytic within $S\times\R^3$ (\lean{AnalyticOn}).
  \item \lean{smoothPeriodicData_analyticOn}: \lean{SmoothPeriodicData}$(u_0)$ implies that $u_0$ is real-analytic.
  \item $((\top:\mathbb{N}_\infty):\leanu{WithTop ℕ∞})\neq\top$.
\end{itemize}
Fefferman's Statement~(B) asks for $C^\infty$ data and $C^\infty$ solutions on
$\R^3\times[0,\infty)$. The repository's \lean{NSOn} is therefore a
\emph{strictly stronger} hypothesis than smoothness, because it requires joint
analyticity in $(t,x)$ up to and including $t=0$ within the time set. For the
results of this paper the issue is harmless in one direction. Every theorem
proved under \lean{NSOn} remains true, and Corollary~\ref{cor:C2} shows the
spatial bridge needs only spatial $C^2$. The issue does, however, affect how the repository's
\lean{FeffermanStatementB}, \lean{PeriodicMaximalSolutionWithBKM} and
\lean{PeriodicBKMIntegrability} compare with their informal counterparts, as
the following remark shows.

\begin{remark}[Informal, \emph{not} formalized: a shear flow that is smooth but not time-analytic at $t=0$]\label{rem:shear}
Fix $\nu>0$. Let
\[
  \varphi_0(r)=\sum_{k\in\mathbb{Z}}e^{-|k|}e^{2\pi ikr}=1+2\sum_{k\ge1}e^{-k}\cos(2\pi kr),
\]
and put $u_0(x)=(0,\varphi_0(x_1),0)$. The series converges for
$|\operatorname{Im}r|<1/(2\pi)$, so $\varphi_0$ is real-analytic on all of
$\R$, and $u_0$ is real-analytic, $1$-periodic and divergence-free, since
$\div u_0=\partial_2\varphi_0(x_1)=0$.

Let $\varphi(t,r)=\sum_ke^{-|k|-4\pi^2\nu k^2t}e^{2\pi ikr}$ for $t\ge0$,
which solves $\varphi_t=\nu\varphi_{rr}$, and let
$u(t,x)=(0,\varphi(t,x_1),0)$ with $p\equiv0$. Then
$(u\cdot\nabla)u=u_2\,\partial_2u=0$, because $u_1=0$ and $u$ does not depend
on $x_2$, and $\Delta u=(0,\varphi_{rr},0)$. So $(u,p)$ is a smooth periodic
solution of the unforced equations on $[0,\infty)\times\R^3$.

Now let $h(t)=\varphi(t,0)$. Every $t$-derivative series is dominated on
$[0,\infty)$ by $\sum_k(4\pi^2\nu k^2)^ne^{-|k|}<\infty$, so the one-sided
derivatives at $0$ are
\[
  h^{(n)}(0)=\sum_k(-4\pi^2\nu k^2)^ne^{-|k|}.
\]
All terms have the sign $(-1)^n$. Keeping only the term $k=2n$ and using
$n!\le n^n$,
\[
  \frac{|h^{(n)}(0)|}{n!}\ \ge\ \frac{(4\pi^2\nu)^n(2n)^{2n}e^{-2n}}{n^n}
  =\Bigl(\frac{16\pi^2\nu\,n}{e^{2}}\Bigr)^{n}.
\]
The $n$-th root tends to $\infty$, so the Taylor series of $h$ at $0$ has
radius of convergence $0$. If $(t,x)\mapsto u(t,x)$ were analytic within
$[0,\infty)\times\R^3$ at $(0,0)$, then so would be its restriction
$t\mapsto u_2(t,0)=h(t)$ within $[0,\infty)$. The coefficients of a power
series that represents $h$ on $[0,r)$ are forced to be $h^{(n)}(0)/n!$, so
that series would have positive radius. This is a contradiction. Hence this
solution does \emph{not} satisfy \lean{NSOn}$(\nu,[0,\infty),u_0,u,p)$,
although $u_0$ satisfies \lean{SmoothPeriodicData}.

Smooth periodic solutions with given data are unique. This is classical, by
the energy method, and it is \emph{not} formalized. Assuming it, any
\lean{NSOn} solution for this $u_0$, which is in particular smooth, would
coincide with the shear flow above, which is impossible. Granting classical
uniqueness, the formal \lean{FeffermanStatementB} is therefore
\emph{false}. The same datum would make the formal
\lean{PeriodicMaximalSolutionWithBKM} false as well, since that statement also
requires \lean{NSOn} on $[0,T)$ with $T>0$, an interval containing $t=0$. The
reduction \lean{feffermanStatementB_of_bkm} remains a correct implication, but
the formal statements do not encode Fefferman's $C^\infty$ problem. None of
this touches the results of this paper. We recommend replacing \leanu{⊤} by
$\infty$, that is \leanu{((⊤ : ℕ∞) : WithTop ℕ∞)}, in \lean{NSOn} and
\lean{SmoothPeriodicData}, in a separate, reviewed change. This paper does not
make that change.
\end{remark}

%==========================================================================
\section{Solution-linked vortex-line packets}\label{sec:packets}

For a velocity $u$ and a time $t$ write $\omega(t,\cdot)=\curl u(t)$ and
$N(t,\cdot)=|\omega(t,\cdot)|$. The normalized vorticity is
\[
  \xi(t,x)=N(t,x)^{-1}\omega(t,x)
\]
(\declref{direction}{Depletion/SolutionLinkedVortexLines.lean}{24}),
with the convention $0^{-1}=0$, so that $\xi=0$ where $\omega=0$. If
$\omega(t,x)\ne0$ then $|\xi(t,x)|=1$
(\declref{norm_direction_of_nonzero}{Depletion/SolutionLinkedVortexLines.lean}{28}).
Its divergence is $\div\xi=\sum_i\partial_i\xi_i$
(\declref{directionDivergence}{Depletion/SolutionLinkedVortexLines.lean}{34}).

\begin{definition}[Packet; \declref{Packet}{Depletion/SolutionLinkedVortexLines.lean}{56}]\label{def:packet}
A \emph{packet} $P$ over $(u,T)$ consists of the following data:
\begin{itemize}
  \item maps $x,y\colon\R\to E$ (the peak and endpoint) and a curve $\gamma\colon\R\times\R\to E$;
  \item a length $L\colon\R\to\R$;
  \item constants $K\in\R_{\ge0}$ with $K\ge1$ and $C\in\R_{\ge0}$.
\end{itemize}
For every $t\in[0,T)$ the data must satisfy:
\begin{enumerate}
  \item[(P1)] $L(t)\ge0$, $\gamma(t,0)=x(t)$ and $\gamma(t,L(t))=y(t)$;
  \item[(P2)] the near-maximum condition $\Omega(t)\le K\,N(t,x(t))$, in $\Ennr$;
  \item[(P3)] if $\omega(t,x(t))\ne0$, then:
    \begin{itemize}
      \item[(a)] $L(t)>0$;
      \item[(b)] $\omega(t,\gamma(t,s))\ne0$ for all $s\in[0,L(t)]$;
      \item[(c)] $s\mapsto\gamma(t,s)$ has derivative $\xi(t,\gamma(t,s))$ within $[0,L(t)]$ at every $s\in[0,L(t)]$, so $\gamma(t,\cdot)$ is a unit-speed parametrization of a vortex line;
      \item[(d)] $s\mapsto(\div\xi)(t,\gamma(t,s))$ is interval-integrable on $[0,L(t)]$;
    \end{itemize}
  \item[(P4)] the line control $\Lambda(t):=\bigl|\int_0^{L(t)}(\div\xi)(t,\gamma(t,s))\,ds\bigr|$ satisfies $\Lambda(t)\le C$ (\lean{lineControl}, line~46; set to $0$ off $[0,T)$).
\end{enumerate}
In addition the following maps must be Borel measurable on all of $\R$: $x$,
$y$, $L$, $\Lambda$, and the endpoint magnitude
\[
  m(t):=\begin{cases}N(t,y(t)) & t\in[0,T),\\ 0 & \text{otherwise}\end{cases}
\]
(\declref{endpointMagnitude}{Depletion/SolutionLinkedVortexLines.lean}{39}).
\end{definition}

Every quantity in (P2)--(P4) and in $m$ is evaluated from the actual field
$\curl u(t)$. There is no free magnitude or direction field that could be
filled in independently of $u$. This is the difference from the legacy free
interface \lean{DHYPacket} (\lean{lean/Clay/Gaps/DirectionAlignment.lean}:28),
which is preserved unchanged. The constant $K$ is the reciprocal of DHY's
near-maximum fraction. Condition (P3a) is an \emph{extra} restriction of this
implementation: DHY do not exclude degenerate segments, and we do not claim
that they do.

\begin{definition}[Endpoint hypothesis and spatial comparison]\label{def:endpoint}
For a packet $P$:
\begin{itemize}
  \item \declref{EndpointIntegralFinite}{Depletion/SolutionLinkedVortexLines.lean}{87} is the statement $\int^{-}_{[0,T)}m(t)\,dt<\infty$;
  \item \declref{SpatialComparison}{Depletion/SolutionLinkedVortexLines.lean}{92} is the statement $N(t,x(t))\le e^{C}\,m(t)$ in $\Ennr$ for all $t\in[0,T)$;
  \item \declref{SpatialBridgeObligation}{Depletion/SolutionLinkedVortexLines.lean}{101} is the statement that \lean{SpatialComparison} holds for every packet over every solution in \lean{ForAllFiniteSolutions}.
\end{itemize}
\end{definition}

\begin{remark}[Nonvacuity, and its limit]
The zero field is a solution (\lean{zero_NSOn}, line~204). It carries the
packet \lean{zeroPacket} (line~164), which satisfies the comparison, the
endpoint hypothesis and BKM finiteness (\lean{zeroPacket_BKMIntegralFinite},
line~229). Conversely, if $L(t)=0$ then $\curl u(t)\equiv0$ everywhere
(\declref{curl_eq_zero_of_length_eq_zero}{Depletion/SolutionLinkedVortexLines.lean}{117}),
so a zero-length placeholder cannot certify a field with nonzero vorticity.
The library contains \emph{no} packet over a nonzero solution. Nonvacuity is
therefore established only for the zero solution.
\end{remark}

%==========================================================================
\section{The spatial bridge: statements and complete proofs}\label{sec:proofs}

The main results are the following two theorems.

\begin{theorem}[\declref{spatialBridgeObligation_proved}{Depletion/ActualSpatialBridge.lean}{173}]\label{thm:bridge}
\lean{SpatialBridgeObligation} holds. More generally
(\declref{spatialComparison_of_NSOn}{Depletion/ActualSpatialBridge.lean}{166}),
if $\mathrm{NSOn}(\nu,[0,T),u_0,u,p)$ holds for some $T\in\Ennr$, which may
be infinite, then every packet $P$ over $(u,T)$ satisfies, for all
$t\in[0,T)$,
\[
  |\omega(t,x(t))|\ \le\ e^{C}\,|\omega(t,y(t))|.
\]
\end{theorem}

\begin{theorem}[\declref{bkmIntegralFinite_of_NSOn_and_endpoint}{Depletion/ActualSpatialBridge.lean}{179}]\label{thm:bkm}
If $\mathrm{NSOn}(\nu,[0,T),u_0,u,p)$ holds and $P$ is a packet over $(u,T)$
satisfying \lean{EndpointIntegralFinite}, then
$\mathrm{BKMIntegralFinite}(u,T)$ holds.
\end{theorem}

The Lean statements, verbatim from the repository, are:
\begin{leanblock}
\kw{theorem} spatialComparison\_of\_NSOn\\
\ \ \ \ \{ν : ℝ\} \{u0 : E3 → E3\} \{u : ℝ → E3 → E3\} \{p : ℝ → E3 → ℝ\} \{T : ℝ≥0∞\}\\
\ \ \ \ (hsol : NSOn ν (lifespanSet T) u0 u p) (P : Packet u T) : SpatialComparison P\\[0.6ex]
\kw{theorem} spatialBridgeObligation\_proved : SpatialBridgeObligation\\[0.6ex]
\kw{theorem} bkmIntegralFinite\_of\_NSOn\_and\_endpoint\\
\ \ \ \ \{ν : ℝ\} \{u0 : E3 → E3\} \{u : ℝ → E3 → E3\} \{p : ℝ → E3 → ℝ\} \{T : ℝ≥0∞\}\\
\ \ \ \ (hsol : NSOn ν (lifespanSet T) u0 u p) (P : Packet u T)\\
\ \ \ \ (hend : EndpointIntegralFinite P) : BKMIntegralFinite u T
\end{leanblock}

The proof is a chain of seven lemmas, each a named Lean theorem. We give every
step.

\subsection{Step 1: spatial \texorpdfstring{$C^2$}{C2} regularity from the solution hypothesis}
\begin{lemma}[\declref{spatialC2_of_NSOn}{Depletion/ActualSpatialBridge.lean}{60}]\label{lem:C2}
If $\mathrm{NSOn}(\nu,S,u_0,u,p)$ and $t\in S$, then each component
$y\mapsto u(t,y)_a$ is $C^2$ on $E$.
\end{lemma}
\begin{proof}
Let $\iota_t\colon E\to\R\times E$, $\iota_t(y)=(t,y)$. This map is the
product of a constant map and the identity, so it is $C^{\top}$ (here $\top=\omega$; see Section~\ref{sec:analytic}); it maps $E$
into $S\times E$ because $t\in S$. By (N1), $U(z)=u(z_1,z_2)$ is $C^{\top}$ on
$S\times E$. The composition rule for \lean{ContDiffOn} then shows that
$U\circ\iota_t=u(t,\cdot)$ is $C^{\top}$ on $\mathrm{univ}$, that is,
$C^{\top}$ on $E$. A map into $E=\leanu{EuclideanSpace ℝ (Fin 3)}$ is
$C^{n}$ if and only if each coordinate is (\lean{contDiff_piLp}). Finally
$2\le\top$ lowers the index (\lean{ContDiff.of_le}). The case $t=0$ needs no
separate treatment, because the composition is taken within $S\times E$.
\end{proof}

\subsection{Step 2: symmetry of mixed partials}
\begin{lemma}[\declref{pd_pd_comm}{Depletion/ActualSpatialBridge.lean}{11}]\label{lem:schwarz}
If $f\colon E\to\R$ is $C^2$, then $\partial_i(\partial_jf)(x)=\partial_j(\partial_if)(x)$ for all $i,j,x$.
\end{lemma}
\begin{proof}
Since $f$ is $C^2$, $Df$ is differentiable at $x$. For a fixed vector $a$, the
map $y\mapsto Df(y)[a]$ is the evaluation of the differentiable map $Df$ at a
constant, so its derivative at $x$ in direction $b$ is $D^2f(x)[b][a]$
(\lean{fderiv_clm_apply}). Hence
$\partial_i(\partial_jf)(x)=D^2f(x)[e_i][e_j]$ and
$\partial_j(\partial_if)(x)=D^2f(x)[e_j][e_i]$. Schwarz's theorem for $C^2$
maps (\lean{ContDiffAt.isSymmSndFDerivAt}) gives
$D^2f(x)[e_i][e_j]=D^2f(x)[e_j][e_i]$.
\end{proof}

\subsection{Step 3: the divergence of a curl vanishes}
\begin{lemma}[\declref{div_curl_eq_zero}{Depletion/ActualSpatialBridge.lean}{24}]\label{lem:divcurl}
If every component of $v\colon E\to E$ is $C^2$, then $\div\curl v\equiv0$.
\end{lemma}
\begin{proof}
Fix $x$. Each $\partial_jv_a$ is differentiable at $x$, because $v_a$ is $C^2$.
The derivative is linear, so for differentiable $f,g$ we have
$\partial_i(f-g)(x)=\partial_if(x)-\partial_ig(x)$. Expanding the definition
of $\curl$ gives
\begin{align*}
\div\curl v(x)
&=\partial_1(\partial_2v_3-\partial_3v_2)+\partial_2(\partial_3v_1-\partial_1v_3)+\partial_3(\partial_1v_2-\partial_2v_1)\\
&=\partial_1\partial_2v_3-\partial_1\partial_3v_2+\partial_2\partial_3v_1-\partial_2\partial_1v_3+\partial_3\partial_1v_2-\partial_3\partial_2v_1 .
\end{align*}
Lemma~\ref{lem:schwarz} applied to $v_3$, $v_2$ and $v_1$ gives
$\partial_1\partial_2v_3=\partial_2\partial_1v_3$,
$\partial_1\partial_3v_2=\partial_3\partial_1v_2$ and
$\partial_2\partial_3v_1=\partial_3\partial_2v_1$. Substituting these and
grouping terms,
\[
(\partial_2\partial_1v_3-\partial_2\partial_1v_3)
+(\partial_3\partial_1v_2-\partial_3\partial_1v_2)
+(\partial_3\partial_2v_1-\partial_3\partial_2v_1)=0 .
\]
In Lean the three substitutions are \lean{pd_pd_comm} at $(0,1)$, $(0,2)$ and
$(1,2)$, and the final cancellation is \lean{ring}. Neither the momentum
equation nor $\div u=0$ is used.
\end{proof}

\begin{lemma}[\declref{curl_differentiable_of_C2}{Depletion/ActualSpatialBridge.lean}{45}]\label{lem:curldiff}
Under the hypothesis of Lemma~\ref{lem:divcurl}, $\curl v$ is differentiable on $E$.
\end{lemma}
\begin{proof}
A map into $E$ is differentiable at $x$ if and only if each coordinate is
(\lean{differentiableAt_piLp}). Each coordinate of $\curl v$ is a difference of
two first partials $\partial_jv_a$, each differentiable at $x$ as in the
previous proof.
\end{proof}

\subsection{Step 4: the product-rule identity for \texorpdfstring{$|\omega|$ and $\xi$}{|omega| and xi}}
\begin{lemma}[\declref{norm_direction_derivative}{Depletion/ActualSpatialBridge.lean}{72}]\label{lem:prod}
Let $w=\curl v$. Suppose that $w$ is differentiable at $x$, that $w(x)\ne0$,
and that $\sum_i\partial_iw_i(x)=0$. Put $N=|w|$ and $\xi=N^{-1}w$. Then
\[
  DN(x)[\xi(x)]\;=\;-N(x)\,(\div\xi)(x).
\]
\end{lemma}
\begin{proof}
\emph{Differentiability.} The Euclidean norm is differentiable away from $0$,
and $w(x)\ne0$, so $N$ is differentiable at $x$ (\lean{DifferentiableAt.norm}).
Since $N(x)\ne0$, $N^{-1}$ is differentiable at $x$, and hence so is the
product $\xi=N^{-1}w$, together with each coordinate $\xi_i$.

\emph{A global factorization.} For every $y\in E$,
\[
  w_i(y)=N(y)\,\xi_i(y).
\]
If $w(y)=0$, both sides are $0$, because $\xi(y)=0^{-1}\cdot0=0$. If
$w(y)\ne0$, then $N(y)\xi_i(y)=N(y)N(y)^{-1}w_i(y)=w_i(y)$. The factorization
holds as an equality of \emph{functions}, which is what allows us to
differentiate it.

\emph{Product rule.} Differentiating at $x$ in the direction $e_i$ gives
\[
  \partial_iw_i(x)=N(x)\,\partial_i\xi_i(x)+\xi_i(x)\,\partial_iN(x).
\]
\emph{Summation.} Summing over $i=1,2,3$ and using the hypothesis $\sum_i\partial_iw_i(x)=0$:
\[
  0=N(x)\sum_i\partial_i\xi_i(x)+\sum_i\xi_i(x)\,DN(x)[e_i]
   =N(x)\,(\div\xi)(x)+DN(x)\Bigl[\sum_i\xi_i(x)e_i\Bigr].
\]
The last equality uses linearity of $DN(x)$. Since $\xi(x)=\sum_i\xi_i(x)e_i$,
this is $0=N(x)(\div\xi)(x)+DN(x)[\xi(x)]$, which rearranges to the claim.
\end{proof}

This is exactly DHY (3.1). Their computation is
$0=\nabla\cdot(|\omega|\xi)=(\xi\cdot\nabla)|\omega|+(\nabla\cdot\xi)|\omega|$,
and the term $DN(x)[\xi(x)]$ is the directional derivative
$(\xi\cdot\nabla)|\omega|$.

\subsection{Step 5: the logarithmic derivative along the vortex line}
\begin{lemma}[\declref{logDerivativeAlong_of_spatialC2}{Depletion/ActualSpatialBridge.lean}{115}]\label{lem:log}
Let $P$ be a packet over $(u,T)$, and suppose that every $u(t)$ with
$t\in[0,T)$ has $C^2$ components. Let $t\in[0,T)$ with
$\omega(t,x(t))\ne0$, and put $L=L(t)$, $\gamma=\gamma(t,\cdot)$. Then
$f(s):=\log N(t,\gamma(s))$ has derivative $-(\div\xi)(t,\gamma(s))$ within
$[0,L]$ at every $s\in[0,L]$. This is the predicate \lean{LogDerivativeAlong}
(line~107).
\end{lemma}
\begin{proof}
Fix $s\in[0,L]$ and write $z=\gamma(s)$. By (P3b), $\omega(t,z)\ne0$. By
Lemma~\ref{lem:curldiff}, $\omega(t,\cdot)$ is differentiable at $z$, so
$N(t,\cdot)$ is differentiable at $z$, as in the proof of
Lemma~\ref{lem:prod}. By (P3c), $\gamma$ has derivative $\xi(t,z)$ within
$[0,L]$ at $s$. The chain rule for a Fr\'echet-differentiable outer map
composed with a within-differentiable curve
(\lean{HasFDerivAt.comp_hasDerivWithinAt}) shows that $s\mapsto N(t,\gamma(s))$
has derivative $DN(t,\cdot)(z)[\xi(t,z)]$ within $[0,L]$ at $s$. By
Lemma~\ref{lem:divcurl}, $\div\omega(t,\cdot)=0$ at $z$. All hypotheses of
Lemma~\ref{lem:prod} therefore hold at $z$, and the derivative equals
$-N(t,z)(\div\xi)(t,z)$. Since $N(t,z)>0$, the chain rule for $\log$
(\lean{HasDerivWithinAt.log}) gives
\[
  f'(s)=\frac{-N(t,z)\,(\div\xi)(t,z)}{N(t,z)}=-(\div\xi)(t,z)
\]
within $[0,L]$. The division is legitimate because $N(t,z)\ne0$; in Lean this
step is \lean{field_simp} followed by \lean{ring}.
\end{proof}

\subsection{Step 6: integration and exponentiation}
\begin{lemma}[\declref{spatialComparison_of_logDerivativeAlong}{Depletion/ActualSpatialBridge.lean}{135}]\label{lem:ftc}
If $P$ satisfies \lean{LogDerivativeAlong}, then \lean{SpatialComparison}$(P)$ holds.
\end{lemma}
\begin{proof}
Fix $t\in[0,T)$ and write $N_x=N(t,x(t))$, $N_y=N(t,y(t))$ and $L=L(t)$.

\emph{Case $N_x=0$.} The left side of the comparison is $0$, and $0\le$
anything in $\Ennr$.

\emph{Case $N_x\ne0$.} Then $\omega(t,x(t))\ne0$, so (P3a) gives $L>0$. By
(P3b) at $s=L$ and (P1), $\omega(t,y(t))=\omega(t,\gamma(t,L))\ne0$, so
$N_y>0$ (\lean{endpoint_nonzero_of_peak_nonzero}, line~127). Let
$f(s)=\log N(t,\gamma(t,s))$ and $g(s)=(\div\xi)(t,\gamma(t,s))$. We check the
hypotheses of the fundamental theorem of calculus in the form
\lean{intervalIntegral.integral_eq_sub_of_hasDerivAt_of_le}:
\begin{itemize}
  \item $0\le L$;
  \item $f$ is continuous on $[0,L]$, because a function with a derivative within a set at a point is continuous within that set at that point;
  \item at each interior point $s\in(0,L)$, the set $[0,L]$ is a neighbourhood of $s$, so the within-derivative $-g(s)$ is a genuine derivative;
  \item $-g$ is interval-integrable on $[0,L]$ by (P3d).
\end{itemize}
Therefore
\[
  \int_0^L(-g(s))\,ds=f(L)-f(0)=\log N_y-\log N_x ,
\]
using $\gamma(t,0)=x(t)$ and $\gamma(t,L)=y(t)$ from (P1). With
$I:=\int_0^Lg(s)\,ds$ and $\int_0^L(-g)=-I$, this rearranges to
\[
  \log N_x=\log N_y+I .
\]
By (P4), $I\le|I|=\Lambda(t)\le C$, hence $\log N_x\le C+\log N_y$. Since
$\exp$ is monotone, $\exp(\log a)=a$ for $a>0$, and
$\exp(C+\log N_y)=e^{C}N_y$, we obtain
\[
  N_x\le e^{C}N_y .
\]
Finally we pass to $\Ennr$. Since $N_x$ and $N_y$ are nonnegative, the
coercions satisfy $\|\cdot\|_+=\mathrm{ofReal}\|\cdot\|$. \lean{ENNReal.ofReal}
is monotone, and $\mathrm{ofReal}(e^{C}N_y)=\mathrm{ofReal}(e^{C})\cdot\mathrm{ofReal}(N_y)$
because $e^{C}\ge0$. Moreover $m(t)=N_y$ for $t\in[0,T)$. This gives the
inequality required by \lean{SpatialComparison}.
\end{proof}

Only the upper bound $N_x\le e^{C}N_y$ is formalized. The reverse inequality
$N_y\le e^{C}N_x$ of DHY (2.2) follows from the same identity, since
$\log N_y=\log N_x-I\le\log N_x+C$, but it is not stated in the library.

\subsection{Step 7: assembly}
\begin{proof}[Proof of Theorem~\ref{thm:bridge}]
By Lemma~\ref{lem:C2} with $S=[0,T)$, every slice $u(t)$ with $t\in[0,T)$ has
$C^2$ components. Lemma~\ref{lem:log} then gives \lean{LogDerivativeAlong}$(P)$,
and Lemma~\ref{lem:ftc} gives \lean{SpatialComparison}$(P)$. This proves
\lean{spatialComparison_of_NSOn}. For \lean{spatialBridgeObligation_proved},
we unfold \lean{ForAllFiniteSolutions}, introduce
$\nu,\,\nu>0,\,u_0,u,p,T$, the data hypothesis, the hypothesis $T\ne\top$ and
$h_{\rm sol}$, and apply the previous statement. The hypotheses $\nu>0$, the
data hypothesis and $T\neq\top$ are not used.
\end{proof}

\begin{proof}[Proof of Theorem~\ref{thm:bkm}]
This is \declref{bkmIntegralFinite_of_spatialComparison_and_endpoint}{Depletion/SolutionLinkedVortexLines.lean}{135}
applied to the comparison from Theorem~\ref{thm:bridge}. Put
$K'=K\cdot\mathrm{ofReal}(e^{C})$. Then $K'<\infty$, because $K\in\R_{\ge0}$
and $e^{C}<\infty$. For $t\in[0,T)$, (P2) and the comparison give
\[
  \Omega(t)\le K\,N(t,x(t))\le K\,e^{C}m(t)=K'm(t).
\]
The function $K'm$ is measurable, because $m$ is measurable by the packet's
measurability field. The monotonicity of the set lower integral,
\lean{setLIntegral_mono}, requires measurability only of the majorant, and it
gives $\int^{-}_{[0,T)}\Omega\le\int^{-}_{[0,T)}K'm$. Measurability of $m$
also lets the constant come out of the integral (\lean{lintegral_const_mul}):
$\int^{-}_{[0,T)}K'm=K'\int^{-}_{[0,T)}m$. The right side is a product of two
finite elements of $\Ennr$, hence finite.
\end{proof}

\begin{corollary}[only spatial $C^2$ is used; \lean{spatialComparison_of_spatialC2} in \lean{PaperChecks.lean}]\label{cor:C2}
Let $u\colon\R\to E\to E$ be any family whose slices $u(t)$, $t\in[0,T)$, have
$C^2$ components. Then every packet over $(u,T)$ satisfies
\lean{SpatialComparison}. No evolution equation, viscosity, incompressibility
of $u$, periodicity or time regularity is needed.
\end{corollary}
\begin{proof}
Compose Lemma~\ref{lem:log} with Lemma~\ref{lem:ftc}. The Lean term is
\lean{spatialComparison_of_logDerivativeAlong P (logDerivativeAlong_of_spatialC2 P hC2)}.
\end{proof}

\begin{remark}
Corollary~\ref{cor:C2} makes precise that the ``Navier--Stokes'' content of
Theorem~\ref{thm:bridge} is only that solutions are spatially $C^2$. The
statement is kinematic: it relates vorticity magnitudes \emph{at one time}
along one vortex line. This agrees with DHY, whose Theorem~1 is stated for
``any 3-D incompressible flow (Euler or Navier--Stokes)'', and whose
Lemma~1 needs only $\omega(\cdot,t)\in C^1$.
\end{remark}

%==========================================================================
\section{What the BKM consequence does and does not give}\label{sec:equiv}

\begin{proposition}[\lean{endpointIntegralFinite_iff_bkm} in \lean{PaperChecks.lean}]\label{prop:equiv}
For every packet $P$ over $(u,T)$, $\mathrm{BKMIntegralFinite}(u,T)$ implies
\lean{EndpointIntegralFinite}$(P)$. If in addition
$\mathrm{NSOn}(\nu,[0,T),u_0,u,p)$ holds, then the two statements are
equivalent.
\end{proposition}
\begin{proof}
For $t\in[0,T)$ we have $m(t)=N(t,y(t))\le\sup_zN(t,z)=\Omega(t)$, and for
$t\notin[0,T)$ we have $m(t)=0\le\Omega(t)$. So $m\le\Omega$ pointwise on all
of $\R$. Monotonicity of the lower Lebesgue integral for the restricted
measure (\lean{lintegral_mono}) needs no measurability, and it gives
$\int^{-}_{[0,T)}m\le\int^{-}_{[0,T)}\Omega<\infty$. The converse direction
is Theorem~\ref{thm:bkm}.
\end{proof}

Proposition~\ref{prop:equiv} is the honest summary of the analytic content.
\emph{Given} a packet, the vortex-line comparison converts the Clay-core
quantity $\int\Omega$ into the endpoint quantity $\int m$, and back, with no
loss. The criterion helps only when one has independent control of the
vorticity at the far endpoint $y(t)$, as in the numerical scenarios DHY
discuss, where $y(t)$ lies outside a collapsing tube. If $y(t)$ is chosen near
the peak, the endpoint hypothesis is no easier than BKM itself. In the
repository's terms, the Clay core \lean{PeriodicBKMIntegrability}
(\lean{lean/Clay/StatementB.lean}:117) remains \emph{open}. Nothing here
reduces it to a weaker or a checkable hypothesis for general data.

%==========================================================================
\section{Ledger: proved, assumed, and open}\label{sec:ledger}

\subsection*{Kernel-checked in the library (commit \texttt{43653aa})}
The following are proved in Lean with standard axioms only:
\begin{itemize}
  \item the ten named theorems of \lean{lean/Clay/Depletion/ActualSpatialBridge.lean} (Table~\ref{tab:decls});
  \item the supporting theorems of \lean{SolutionLinkedVortexLines.lean} listed in the table;
  \item the conditional reduction \lean{feffermanStatementB_of_bkm} (\lean{lean/Clay/StatementB.lean}:122), from the earlier monograph~\cite{ClayMonograph}.
\end{itemize}

\subsection*{Kernel-checked in this paper's check file (not part of the library audit)}
\begin{itemize}
  \item \lean{NSOn_velocity_analyticOn} and \lean{smoothPeriodicData_analyticOn} (Section~\ref{sec:analytic});
  \item \lean{spatialComparison_of_spatialC2} (Corollary~\ref{cor:C2});
  \item \lean{endpointIntegralFinite_of_bkm} and \lean{endpointIntegralFinite_iff_bkm} (Proposition~\ref{prop:equiv}).
\end{itemize}
The pull request's CI runs this file and prints \lean{#print axioms} for each declaration.

\subsection*{Hypotheses carried by the theorems (assumed, not proved)}
\begin{enumerate}
  \item \textbf{All of \lean{NSOn}.} This means (N1)--(N5), including joint analyticity (Section~\ref{sec:analytic}). Only the spatial $C^2$ consequence of (N1) is used.
  \item \textbf{All packet fields (P1)--(P4)} and the five measurability fields, for the specific packet supplied: near-maximality with a fixed $K$, nonvanishing along the segment, the unit-speed vortex-line tangent equation, integrability of $\div\xi$ along the segment, and the time-uniform bound $C$.
  \item \textbf{\lean{EndpointIntegralFinite}} (Theorem~\ref{thm:bkm} only), a separate argument which by Proposition~\ref{prop:equiv} is as strong as BKM finiteness.
\end{enumerate}

\subsection*{Classical inputs \emph{not} formalized}
\begin{enumerate}
  \item Local existence and uniqueness of smooth periodic solutions, and the BKM continuation criterion~\cite{BKM1984}. These appear only as the named proposition \lean{PeriodicMaximalSolutionWithBKM} (\lean{lean/Clay/StatementB.lean}:109). The step from BKM finiteness to ``no blow-up'' therefore has no formal proof here. In particular Theorem~\ref{thm:bkm} concludes integral finiteness, \emph{not} regularity.
  \item Existence of vortex-line packets for nonzero solutions: integral curves of $\xi$ through $x(t)$, their nonvanishing extent, a measurable choice of $x(t),y(t),L(t)$ in $t$, and time-uniform constants $K$ and $C$. None of this is proved for any nonzero solution.
  \item The general and degenerate configurations of DHY Theorem~1 (for example $L(t)=0$ with nonzero vorticity), and the two-sided ratio bound (2.2).
  \item DHY Theorem~2 and Lemmas~2--3. These use the Euler vorticity transport $\omega(X(\alpha,t),t)=\nabla_\alpha X\,\omega(\alpha,t_1)$, which fails for Navier--Stokes. None of them is imported.
\end{enumerate}

\subsection*{Open}
The Clay core \lean{PeriodicBKMIntegrability} and Fefferman's Statement~(B)
remain open. The fidelity issue of Section~\ref{sec:analytic} is unresolved:
the formal \lean{FeffermanStatementB} requires analytic solutions, and by
Remark~\ref{rem:shear} (informal, and conditional on classical uniqueness) it
is false as written.

%==========================================================================
\section{Declaration index}\label{sec:index}
Locations are given relative to \lean{lean/Clay/} at commit \texttt{\mainsha}.
The CI script \lean{papers/spatial-bridge/check_refs.py} checks every
\emph{file:line} pair printed in this paper: it confirms that the cited line
declares the cited name.

{\small
\begin{longtable}{@{}p{0.47\textwidth}p{0.45\textwidth}@{}}
\caption{Lean declarations used in this paper.}\label{tab:decls}\\
\toprule
Declaration & Location \\ \midrule
\endfirsthead
\toprule Declaration & Location \\ \midrule \endhead
\multicolumn{2}{@{}l}{\emph{New in pull request \#9 (\lean{ActualSpatialBridge.lean}); ten named theorems}}\\
\declref{pd_pd_comm}{Depletion/ActualSpatialBridge.lean}{11} & Lemma~\ref{lem:schwarz} \\
\declref{div_curl_eq_zero}{Depletion/ActualSpatialBridge.lean}{24} & Lemma~\ref{lem:divcurl} \\
\declref{curl_differentiable_of_C2}{Depletion/ActualSpatialBridge.lean}{45} & Lemma~\ref{lem:curldiff} \\
\declref{spatialC2_of_NSOn}{Depletion/ActualSpatialBridge.lean}{60} & Lemma~\ref{lem:C2} \\
\declref{norm_direction_derivative}{Depletion/ActualSpatialBridge.lean}{72} & Lemma~\ref{lem:prod} \\
\declref{LogDerivativeAlong}{Depletion/ActualSpatialBridge.lean}{107} & definition (Prop), Lemma~\ref{lem:log} \\
\declref{logDerivativeAlong_of_spatialC2}{Depletion/ActualSpatialBridge.lean}{115} & Lemma~\ref{lem:log} \\
\declref{spatialComparison_of_logDerivativeAlong}{Depletion/ActualSpatialBridge.lean}{135} & Lemma~\ref{lem:ftc} \\
\declref{spatialComparison_of_NSOn}{Depletion/ActualSpatialBridge.lean}{166} & Theorem~\ref{thm:bridge} \\
\declref{spatialBridgeObligation_proved}{Depletion/ActualSpatialBridge.lean}{173} & Theorem~\ref{thm:bridge} \\
\declref{bkmIntegralFinite_of_NSOn_and_endpoint}{Depletion/ActualSpatialBridge.lean}{179} & Theorem~\ref{thm:bkm} \\ \midrule
\multicolumn{2}{@{}l}{\emph{From pull request \#8 (\lean{SolutionLinkedVortexLines.lean})}}\\
\declref{direction}{Depletion/SolutionLinkedVortexLines.lean}{24} & $\xi$ \\
\declref{norm_direction_of_nonzero}{Depletion/SolutionLinkedVortexLines.lean}{28} & $|\xi|=1$ \\
\declref{directionDivergence}{Depletion/SolutionLinkedVortexLines.lean}{34} & $\div\xi$ \\
\declref{endpointMagnitude}{Depletion/SolutionLinkedVortexLines.lean}{39} & $m$ \\
\declref{lineControl}{Depletion/SolutionLinkedVortexLines.lean}{46} & $\Lambda$ \\
\declref{Packet}{Depletion/SolutionLinkedVortexLines.lean}{56} & Definition~\ref{def:packet} \\
\declref{EndpointIntegralFinite}{Depletion/SolutionLinkedVortexLines.lean}{87} & Definition~\ref{def:endpoint} \\
\declref{SpatialComparison}{Depletion/SolutionLinkedVortexLines.lean}{92} & Definition~\ref{def:endpoint} \\
\declref{SpatialBridgeObligation}{Depletion/SolutionLinkedVortexLines.lean}{101} & Definition~\ref{def:endpoint} \\
\declref{curl_eq_zero_of_length_eq_zero}{Depletion/SolutionLinkedVortexLines.lean}{117} & zero-length rejection \\
\declref{endpoint_nonzero_of_peak_nonzero}{Depletion/SolutionLinkedVortexLines.lean}{127} & Lemma~\ref{lem:ftc} \\
\declref{bkmIntegralFinite_of_spatialComparison_and_endpoint}{Depletion/SolutionLinkedVortexLines.lean}{135} & Theorem~\ref{thm:bkm} \\
\declref{zeroPacket}{Depletion/SolutionLinkedVortexLines.lean}{164} & nonvacuity (zero field) \\
\declref{zero_NSOn}{Depletion/SolutionLinkedVortexLines.lean}{204} & nonvacuity (zero field) \\
\declref{zeroPacket_BKMIntegralFinite}{Depletion/SolutionLinkedVortexLines.lean}{229} & nonvacuity (zero field) \\ \midrule
\multicolumn{2}{@{}l}{\emph{Base definitions (\lean{StatementB.lean})}}\\
\declref{pd}{StatementB.lean}{41} & $\partial_j$ \\
\declref{DivFree}{StatementB.lean}{50} & $\div=0$ \\
\declref{curl}{StatementB.lean}{53} & $\curl$ \\
\declref{SmoothPeriodicData}{StatementB.lean}{60} & data class (Section~\ref{sec:analytic}) \\
\declref{NSOn}{StatementB.lean}{65} & solution predicate \\
\declref{FeffermanStatementB}{StatementB.lean}{77} & open target \\
\declref{lifespanSet}{StatementB.lean}{82} & $[0,T)$ \\
\declref{vortSup}{StatementB.lean}{95} & $\Omega$ \\
\declref{BKMIntegralFinite}{StatementB.lean}{98} & BKM integral \\
\declref{ForAllFiniteSolutions}{StatementB.lean}{102} & quantifier \\
\declref{PeriodicMaximalSolutionWithBKM}{StatementB.lean}{109} & classical, \emph{not formalized} \\
\declref{PeriodicBKMIntegrability}{StatementB.lean}{117} & Clay core, \emph{open} \\
\declref{feffermanStatementB_of_bkm}{StatementB.lean}{122} & proved reduction \\
\declref{DHYPacket}{Gaps/DirectionAlignment.lean}{28} & legacy free interface \\
\bottomrule
\end{longtable}}

\paragraph{Axiom audit.} CI for the merge commit \texttt{43653aa} ran the
library audit \lean{lean/AxiomReport.lean}, which iterates over every
non-internal constant in the \lean{Clay} module tree and collects its axioms.
It reported \texttt{AXSUMMARY 125 theorems, 236 declarations in Clay;
nonstandard: 0}, and every row lists a subset of \texttt{propext},
\texttt{Classical.choice} and \texttt{Quot.sound}. The workflow fails on any
occurrence of \texttt{sorryAx} and on any nonstandard axiom. The count of 125
includes compiler-generated equation and structure lemmas. Twelve of the 125
rows are attributed to the new module: the ten named theorems above, the
auxiliary \lean{LogDerivativeAlong.proof_1}, and the equation lemma
\lean{Clay.StatementB.e.eq_1}, which Lean generates on demand in that
module.

%==========================================================================
\section{Reproducibility}\label{sec:repro}
\begin{itemize}
  \item \textbf{Repository:} \url{https://github.com/BenFrohman/Clay-}.
  \item \textbf{Library commit:} \texttt{\mainsha} (\texttt{main}; merge of pull request~\#9). The Lean files cited here are unchanged on the paper branch.
  \item \textbf{Toolchain:} \texttt{leanprover/lean4:v4.16.0} (\lean{lean/lean-toolchain}).
  \item \textbf{Mathlib:} \texttt{v4.16.0}, revision \texttt{a6276f4c6097675b1cf5ebd49b1146b735f38c02}, pinned in \lean{lean/lake-manifest.json} together with its transitive dependencies.
  \item \textbf{CI:} the workflow \lean{.github/workflows/lean.yml} (Lean build, axiom audit, regressions and paper checks) and \lean{.github/workflows/paper-spatial-bridge.yml} (reference check and LaTeX build).
\end{itemize}
To reproduce, run:
\begin{lstlisting}
git clone https://github.com/BenFrohman/Clay- && cd Clay-
git checkout 43653aab3b51269a71a713ccbc4d892869bbe497   # or the paper branch
cd lean
lake exe cache get          # prebuilt Mathlib for the pinned revision
lake build                  # builds the Clay library
lake env lean AxiomReport.lean | grep AXSUMMARY
#   expected: AXSUMMARY 125 theorems, 236 declarations in Clay; nonstandard: 0
lake env lean SolutionLinkedRegression.lean
lake env lean ../papers/spatial-bridge/PaperChecks.lean   # paper branch only
cd ../papers/spatial-bridge
python3 check_refs.py       # every file:line in the paper names its declaration
tectonic spatial_bridge.tex # builds spatial_bridge.pdf
\end{lstlisting}
In CI the clean library build takes about four minutes with the Mathlib cache.

\paragraph{Acknowledgement of method.} The Lean development and this paper
were prepared with AI-assisted tooling under the author's direction. The
author is responsible for the content. All formal claims are those checked by
the Lean kernel as described above.

\bibliographystyle{plain}
\bibliography{spatial_bridge}

\end{document}
